\documentclass[12pt,a4paper,oneside]{report}

\usepackage[utf8]{inputenc}
\usepackage[T1]{fontenc}
\usepackage[italian,english]{babel}
\usepackage{geometry}
\usepackage{setspace}
\usepackage{titlesec}
\usepackage{graphicx}
\usepackage{emptypage}
\usepackage{hyperref}

\geometry{
    top=3cm,
    bottom=3cm,
    left=3cm,
    right=3cm
}

\setstretch{1.2}

\begin{document}

\begin{titlepage}
    \begin{center}
        \vspace*{1cm}

        {\Large \textbf{UNIVERSIT\`A DEGLI STUDI DI TRENTO}}\\[0.4cm]
        {\large Dipartimento di Ingegneria e Scienza dell'Informazione}\\[1.5cm]

        {\large Corso di Laurea in Informatica}\\[2cm]

        {\large \textbf{Tesi di Laurea}}\\[1.8cm]

        {\Huge \textbf{Quantum Annealing Learning Search}}\\[0.4cm]
        {\Large \textbf{per la risoluzione di problemi QUBO}}\\[0.8cm]

        {\large
        Studio dell'algoritmo QALS su problemi QUBO
        }\\[2.5cm]

       \noindent
        \begin{minipage}[t]{0.45\textwidth}
            \vspace{0pt}
            \textbf{Laureando:}\\
            Francesco Buscardo\\
            \end{minipage}
            \hfill
            \begin{minipage}[t]{0.45\textwidth}
            \vspace{0pt}
            \raggedleft
            \textbf{Relatore:}\\
            Prof.\ Enrico Blanzieri
        \end{minipage}
        \vfill

        {\large Anno Accademico 2025/2026}
    \end{center}
\end{titlepage}

\clearpage
\thispagestyle{empty}
\vspace*{\fill}

\begin{center}
\begin{minipage}{0.5\textwidth}
\begin{ttfamily}
\raggedright
while\ (!task\_is\_completed) \{\\
\hspace*{1.5em}try\_solve\_problem();\\
\}
\end{ttfamily}
\end{minipage}
\end{center}

\vspace*{\fill}
\clearpage

%========================
% ABSTRACT
%========================
\chapter*{Abstract}
\addcontentsline{toc}{chapter}{Abstract}

L'obiettivo del lavoro è esplorare come \emph{Quantum Annealing Learning Search} (QALS) riesca
a risolvere problemi di \emph{ottimizzazione quadratica binaria} (QUBO) e confrontare le sue prestazioni
con altri metodi classici.

In particolare i problemi su cui si concentrerà lo studio sarà la seguente varianti del knapsack problem con soft constraints:
\emph{{0/1} Knapsack Problem}.

\clearpage

%========================
% RINGRAZIAMENTI
%========================
\chapter*{Ringraziamenti}
\addcontentsline{toc}{chapter}{Ringraziamenti}

Desidero ringraziare il mio relatore per il supporto, la disponibilit\`a e i consigli
forniti durante lo sviluppo di questo lavoro.

Un ringraziamento speciale va alla mia famiglia e alle persone a me vicine, che mi
hanno sostenuto nel corso di tutto il percorso universitario.

Ringrazio inoltre i colleghi e gli amici con cui ho condiviso studio, dubbi e momenti
difficili, contribuendo a rendere questo percorso pi\`u leggero e significativo.

\clearpage

%========================   
% INDICE
%========================
\tableofcontents

\clearpage

%========================
% CAPITOLI
%========================

\chapter{Introduzione}
\label{ch:introduzione}

\section{Motivazioni}
\label{sec:motivazioni}
% Spiegare perché studiare QALS e il quantum annealing applicato a problemi QUBO;
% contesto dell'informatica quantistica e limiti dei metodi classici che motivano l'approccio ibrido.

\section{Obiettivi}
\label{sec:obiettivi}
% Definire cosa si vuole dimostrare/verificare nella tesi: analizzare le prestazioni di QALS
% sulle varianti del Knapsack Problem, confrontarle con metodi classici e proporre miglioramenti.

\chapter{Background teorico}
\label{ch:background}

\section{Quantum Annealing}
\label{sec:quantum-annealing}
% Descrivere il principio fisico del quantum annealing, l'evoluzione adiabatica verso lo stato
% fondamentale, l'hardware D-Wave e le topologie (Chimera, Pegasus) con i relativi limiti di embedding.

\chapter{Quantum Annealing Learning Search (QALS)}
\label{ch:qals}

\section{Algoritmo}
\label{sec:qals-algoritmo}
% Descrivere l'idea generale e lo pseudocodice di QALS: ciclo di annealing classico esterno,
% meccanismo tabù integrato nel profilo energetico, notazione e parametri principali.

\section{Dimostrazione di convergenza}
\label{sec:qals-convergenza}
% Riportare le ipotesi del teorema di convergenza per catene di Markov disomogenee
% e la verifica di tali ipotesi per QALS, come mostrato nel paper di Pastorello e Blanzieri (2019).

\section{Implementazioni esistenti}
\label{sec:qals-implementazioni}
% Riassumere le implementazioni C++ e Python di QALS descritte in Bonomi et al. (QI\&C 2022),
% le scelte implementative e i risultati ottenuti su Number Partitioning Problem e Travelling Salesman Problem.

\chapter{Knapsack Problem}
\label{ch:knapsack}

\section{Formulazione QUBO}
\label{sec:knapsack-qubo}
% Definire le varianti del Knapsack Problem trattate (0/1, Multidimensional, Quadratic)
% e introdurre il termine di penalità quadratico per trasformarle in problemi QUBO.

\section{Derivazione della matrice Q}
\label{sec:knapsack-matrice-q}
% Sviluppare passo per passo il calcolo algebrico che porta dalla funzione obiettivo penalizzata
% alla matrice $Q$, distinguendo i termini lineari e quadratici.

\section{Scelte del parametro $\lambda$}
\label{sec:knapsack-lambda}
% Confrontare i tre metodi di calcolo di $\lambda$ utilizzati nei report ($\sum p_i / C$,
% $\sum p_i / 3$, $\sum p_i / (650 \cdot C)$) e il loro effetto sulla penalizzazione del vincolo.

\chapter{Metodologia sperimentale}
\label{ch:metodologia}

\section{Baseline classica}
\label{sec:baseline-classica}
% Descrivere il confronto tra QALS e un approccio classico (non quantistico) sulle stesse istanze,
% con le metriche Weight GAP e Profit GAP.

\section{Scaling QUBO}
\label{sec:scaling-qubo}
% Spiegare l'effetto della moltiplicazione della matrice $Q$ per un fattore $\alpha$ sull'energia
% e sulla probabilità di accettazione nel simulated annealing, con il compromesso tra esplorazione e precisione.

\section{Parametro Swap(k, TIMES)}
\label{sec:swap-k-times}
% Introdurre la distinzione tra il parametro $k$ (numero di chiamate all'annealer) e TIMES
% (numero di ripetizioni di QALS/annealer), motivandone l'introduzione e i risultati preliminari.

\section{Confronto con Gurobi}
\label{sec:confronto-gurobi}
% Presentare la validazione delle soluzioni trovate da QALS confrontandole con quelle ottenute da Gurobi
% sulle istanze di test, evidenziando corrispondenze e discrepanze.

\chapter{Diagnostica di QALS}
\label{ch:diagnostica}

\section{Distanze di Hamming}
\label{sec:distanze-hamming}
% Definire le quattro metriche di distanza di Hamming usate per analizzare il comportamento di QALS:
% tra soluzione proposta e corrente, proposta e ottima, soluzioni accettate consecutive, accettata e ottima.

\section{Clustering agglomerativo}
\label{sec:clustering}
% Descrivere la costruzione della matrice delle distanze, il clustering agglomerativo delle soluzioni
% esplorate, l'individuazione del medoid e l'interpretazione in termini di concentrazione/dispersione della ricerca.

\section{Decay factor}
\label{sec:decay-factor}
% Presentare l'introduzione del decay factor nella matrice tabù, la sua formulazione matematica
% e gli effetti osservati sui risultati sperimentali.

\chapter{Proposta di miglioramento}
\label{ch:proposta}
% Illustrare le criticità individuate nella diagnostica (es. concentrazione della ricerca in poche regioni),
% la modifica proposta all'algoritmo, la sua motivazione teorica e il confronto con la versione originale.

\chapter{Risultati sperimentali}
\label{ch:risultati}

\section{Riepilogo comparativo per istanza}
\label{sec:riepilogo-istanze}
% Riassumere e confrontare i risultati ottenuti su tutte le istanze di test (QALS vs. classico vs. Gurobi,
% effetto dei parametri $\lambda$, $k$/TIMES, decay factor), con discussione critica.

\chapter{Conclusioni}
\label{ch:conclusioni}
% Riassumere il lavoro svolto, i risultati principali, i limiti dello studio e i possibili sviluppi futuri.

\appendix
\chapter{Appendici}
\label{ch:appendici}

\section{Pseudocodice}
\label{sec:pseudocodice}
% Riportare per esteso lo pseudocodice di QALS (versione originale e versione modificata).

\section{Dati delle istanze}
\label{sec:dati-istanze}
% Elencare i dati completi delle istanze di test utilizzate (ksp\_1, ksp\_2, ksp\_3).

\section{Tabelle}
\label{sec:tabelle}
% Includere le tabelle riassuntive complete dei risultati sperimentali non riportate nel corpo del testo.

%========================
% BIBLIOGRAFIA
%========================
\begin{thebibliography}{9}

\bibitem{pastorello2019}
D. Pastorello, E. Blanzieri,
\textit{Quantum annealing learning search for solving QUBO problems},
Quantum Information Processing, 2019.

\bibitem{bonomi2022}
A. Bonomi, T. De Min, E. Zardini, E. Blanzieri, V. Cavecchia, D. Pastorello,
\textit{Quantum Annealing Learning Search Implementations},
Quantum Information and Computation, 2022.


\bibitem{demin2021qalscpp}
T.~De Min,
\newblock \emph{C++ Quantum Annealing Learning Search},
\newblock 2021,
\newblock \url{https://github.com/tdemin16/QALS-cpp}.

\bibitem{bonomi2021qalspython}
A.~Bonomi,
\newblock \emph{Python Quantum Annealing Learning Search},
\newblock 2021,
\newblock \url{https://github.com/bonom/Quantum-Annealing-for-solving-QUBO-Problems}.

\bibitem{wikipedia_multiple_constraints}
Wikipedia contributors,
\newblock \emph{List of knapsack problems --- Multiple constraints},
\newblock Wikipedia,
\newblock \url{https://en.wikipedia.org/wiki/List_of_knapsack_problems#Multiple_constraints}.

\bibitem{wikipedia_quadratic_knapsack}
Wikipedia contributors,
\newblock \emph{Quadratic knapsack problem},
\newblock Wikipedia,
\newblock \url{https://en.wikipedia.org/wiki/Quadratic_knapsack_problem}.

\bibitem{ieee_qals}
Evren G\"{u}ney, Joachim Ehrenthal, Thomas Hanne,
\newblock \emph{QUBO Formulations and Characterization of Penalty Parameters for the Multi-Knapsack Problem},
\newblock \url{https://ieeexplore.ieee.org/abstract/document/10924197}

\bibitem{pymhlib}
\newblock \emph{pymhlib --- Python Metaheuristics Library},
\newblock \url{https://github.com/ac-tuwien/pymhlib.git}.

\bibitem{pysolution}
\newblock \emph{pysolution --- Python Knapsack Solution},
\newblock \url{https://github.com/mohitsingla123/Dynamic-Programming-java/blob/master/0-1%20knapsack/pysolution.py}

\bibitem{pisinger_generator}
D. Pisinger,
\newblock \emph{Generator for the Knapsack Problem},
\newblock \url{https://hjemmesider.diku.dk/~pisinger/generator.c}

\end{thebibliography}

%========================
% ULTIMA PAGINA
%========================
\clearpage
\thispagestyle{empty}
\vspace*{\fill}
\begin{center}
    \textit{Fine}
\end{center}
\vspace*{\fill}

\end{document}